\section{Limitations and Conclusion}
\label{sec:limitations}

\paragraph{The models are fitted on a moving target.}
Mean lead time falls 5.34 days between training and test, and a model fitted on
the slower past is systematically off-level on the faster present ---
Section~\ref{sec:results-segments} measures that as a positive bias in every
region. Hence $R^2$ against a test-period variance is negative even for a useful
model, which is why MAE leads, and hence the models need refitting on a rolling
basis rather than once \citep{gama2014drift}; calibration belongs in the same
cycle, because a calibration fitted on one prior drifts as the prior does.

\paragraph{The largest driver of lateness is not observable at prediction time.}
The late rate swings between 1.4\% and 21.4\% month to month, and
Section~\ref{sec:eda-lateness} traces the swings to a Black Friday surge, a
carrier-side slowdown that coincided with a tightened promise, and a national
strike absorbed into a longer promised window: shocks to carrier capacity or to
the platform's own promising policy, none an attribute of an order. This, not
the choice of algorithm, is the binding constraint on Problem 1B, and it is why
the ceiling was stated in advance rather than a weak classifier presented as a
disappointment.

\paragraph{The benefit of outreach is assumed, not measured.}
Section~\ref{sec:cost} values a contact at the review it prevents, and the data
contain no contacts: nobody in this history was reached before reviewing, so the
effect of reaching them --- and the possibility that an unsolicited contact
lowers a rating --- is unobserved. The precision and recall figures stand on
their own; the reais figures do not, and should be read as the value of the
policy \emph{if} the pilot of Section~\ref{sec:pilot} finds the contact averts at
least one low review in five. Part of the assumption can be checked, and fails
for late parcels and for orders that never arrive: most of their low reviews
are written before any delivery, so the value is restated on the reachable share
(Appendix~\ref{app:timing}). Until the
pilot runs, the recommendation is to pilot for six weeks, not to scale.

\paragraph{What the data does not contain, and what rests on assumption.}
The target of Problem 2 is a platform rating, with the self-selection
Section~\ref{sec:problem} describes: a model of the rating is not a model of
satisfaction, and the 1--2 star label cannot see the quietly dissatisfied
customer who never reviews. The service-quality literature identifies product
expectation gaps, seller communication and returns handling among the drivers of
dissatisfaction \citep{parasuraman2005esqual, rao2011elsq}; none appears in these
nine tables, so the Problem 2 residual is unobserved cause rather than noise, and
the free-text review columns are too sparse to recover it.
Selection votes on three out-of-sample windows --- the right trade under drift,
but a small sample: Section~\ref{sec:results-late} records a case where the
rule preferred a model that does not transfer, and
Section~\ref{sec:results-review} shows a fold-to-fold spread too wide to
separate the best candidates at all.

\paragraph{Generalisation, and what these scores are not.}
Everything here is Brazil, 2017--2018, on one marketplace whose geography ---
59.7\% of sellers in S\~{a}o Paulo state against 42.0\% of customers, a median
order travelling 434\,km --- the models' distance-driven structure reflects.
The outputs are triage scores for internal prioritisation, not
entitlements, and must not be used to penalise a seller, price a customer, or
drive an automated decision about someone without review; the review-risk model
is largely a restatement of lateness, so a seller penalised by it would often
be penalised for a carrier's failure.

\subsection*{Conclusion}

Two problems, at two prediction moments, from one dataset and one feature
pipeline, carried through to a deployed application. The regression reaches
MAE 3.90 days --- 41\% below a constant fitted on the past and 60\% below the
platform's own promise --- and its output, thresholded against that promise,
makes a better late-order flag (2.89$\times$ the no-skill floor) than any
classifier trained directly on the binary label (1.24$\times$), so the brief's
dual framing turns out to be the substantively better design. The satisfaction
model reaches 3.75$\times$ its floor on a target where always predicting
``satisfied'' scores 89\% accuracy and catches nothing, and read as a ranked
list it turns ten contacts a day into an outreach queue of 57\% precision ---
a value limited by capacity, not by the model's threshold.

The methodological findings travel further. Cross-validation inside a drifting
training period reorders the candidates relative to the forward window, and its
fold-to-fold spread is too wide to tell the best candidates apart --- so
selection belongs on several recent out-of-sample windows judged by rank, and
even then it can fail. Weighting the loss for class imbalance leaves ranking
unchanged and inflates predicted risk threefold, so the threshold is the right
lever. And the most useful business results are
negative: an informative flag that does not justify an expensive intervention,
and a paper gain of R\$775--1{,}466 per thousand orders that is worth a six-week
experiment before it is worth anything.
